\chapter{Mathematical Notation and Glossary}
\label{apx:notation}

To ensure formal rigor and clarity of exposition, the subsequent chapters utilize a mathematical notation derived from type theory and functional semantics. This chapter serves as a quick reference for the symbology and formal concepts introduced in the formalization of the system.

\section{Notational Conventions}
The model uses the following standard conventions to describe the data structures and functions of the system:

\begin{itemize}
    \item \textbf{Tuples (Cartesian Product)} $\langle x_1, x_2, \dots, x_n \rangle$: Represent composite types (corresponding to \texttt{structs} in the source code). Each element of the tuple defines an essential property of the entity.
    \item \textbf{Disjoint Union (Sum Type)} $A \mid B \mid C$: Indicates that a value can belong to one and only one type among those listed (corresponding to \texttt{enums} in the code). It is used to model polymorphism or uncertainty.
    \item \textbf{Vectors and Sequences} $\vec{X}$: Indicates a finite, ordered sequence of elements of type $X$ (corresponding to lists, arrays, or \texttt{Vec} in memory).
    \item \textbf{Optionality} $X^?$: Indicates a value that may be present (of type $X$) or absent (corresponding to \texttt{Option<T>}).
    \item \textbf{Function Signatures} $f : A \to B$: Denotes a function $f$ that takes an input from domain $A$ and returns an output in codomain $B$. If the input consists of multiple arguments, the Cartesian product is used (e.g., $A \times B \to C$).
\end{itemize}

\section{Glossary of Types and Domains}
The following defines the principal symbols that constitute the alphabet of our Intermediate Representation (IR) and the analysis pipeline.

\subsection{Base Types and Identifiers}
\begin{itemize}
    \item $\mathcal{I}$ (\textbf{Identifier}): The set of all valid alphanumeric strings used to name variables, functions, or classes.
    \item $\pi$ (\textbf{Qualified Name}): A hierarchical sequence of identifiers ($\langle \iota_1, \dots, \iota_n \rangle \in \mathcal{I}^*$) that represents the absolute and unique path of an entity (e.g., a namespace or a module path).
    \item $\tau$ (\textbf{Type Reference}): An abstract reference to a type. It is represented by a seven-state sum type that models its lifecycle during the analysis: \texttt{Prim} (base primitive), \texttt{Unresolved} (raw name just extracted), \texttt{ResolutionQuery} (algebraic expression of intent), \texttt{Resolved} (verified absolute path in the codebase), \texttt{External} (path resolved via import but not analyzed in the codebase), \texttt{Failed} (unsuccessful resolution attempt), or \texttt{EvaluatedAccess} (a wrapper that preserves both the originally queried path and the final structural type, ensuring the generation of multiple edges for intermediate access chains).
\end{itemize}

\subsection{Structural Members and Signatures}
\begin{itemize}
    \item $v$ (\textbf{Variable}): A local variable declared within a block, defined by its identifier and type ($\langle \iota, \tau \rangle$).
    \item $f$ (\textbf{Field}): A data field belonging to a structure, defined by the identifier and type pair ($\langle \iota, \tau \rangle$).
    \item $p$ (\textbf{Parameter}): An input parameter to a function, comprising an optional name, a type, and a variadicity indicator ($\langle \iota^?, \tau, \mathbb{B} \rangle$).
    \item $\sigma$ (\textbf{Signature}): The architectural signature of a function, composed of the tuple $\langle \vec{p}, \tau_{ret} \rangle$. It abstracts the input and output types, isolating them from the implementing component.
    \item $\mathcal{B}$ (\textbf{Block}): A lexical block (e.g., delimited by \texttt{\{\}}) that models an internal scope. It tracks declared local variables, calls, instantiations, and recursively nests any sub-blocks.
    \item $\mathcal{F}$ (\textbf{Function}): A universal functional component defined by the tuple $\langle \iota, \sigma, \mathcal{B} \rangle$. Its architectural role (method or free function) is determined solely by the construct in which it is nested.
\end{itemize}

\subsection{Architectural Entities}
\begin{itemize}
    \item $\mathcal{S}$ (\textbf{Structured Type}): The abstraction of classes, structs, interfaces, and traits. It aggregates fields, functions, inheritance hierarchies, and nested types.
    \item $\mathcal{A}$ (\textbf{Alias}): A type alias mapping a new identifier to an existing type reference ($\langle \iota, \tau \rangle$), such as a \texttt{typedef} in C or a \texttt{type} alias in Rust.
    \item $\mathcal{X}$ (\textbf{Implementation Extension}): A temporary syntactic extraction construct typical of languages like Rust or C++ (the so-called \textit{Impl Blocks}). It associates a set of functions and aliases to a target type. It is flattened into $\mathcal{S}$ during the Resolution Phase, disappearing from the final model.
    \item $\mathcal{C}$ (\textbf{Component}): The universal sum type that embraces the definitive analyzable entities of the system (Modules, Structured Types, Functions, and Aliases).
    \item $\mathcal{M}$ (\textbf{Module}): The fundamental hierarchical unit. The root module acts as the domain of the Intermediate Representation ($\mathcal{D}_{IR}$) and recursively aggregates all analyzed source code.
\end{itemize}

\subsection{Resolution Model and Graph}
\begin{itemize}
    \item $\mathcal{Q}$ (\textbf{Query Algebra}): The set of constant algebraic expressions (\text{Find}, \text{Extract}, \text{Call}) that model the logical intent of resolving a type, deferring its execution (lazy evaluation).
    \item $\Sigma$ (\textbf{Scope}): A single environment of lexical visibility, containing a map of symbols associated with identifiers $\iota$, and references to the base types (supertypes) from which it inherits.
    \item $\mathcal{T}$ (\textbf{Scope Tree}): The global tree of $\Sigma$ scopes. It constitutes the central Name Resolution architecture of the system. It replaces classic flat registries, allowing hierarchical resolution (lexical climbing) and inheritance in logarithmic time with respect to the height of the tree.
    \item $\Gamma$ (\textbf{Execution Context}): The immutable context composed of the tree $\mathcal{T}$, the ecosystem of primitive types, and the current base path, used by the evaluator to execute the queries $\mathcal{Q}$.
    \item $\mathcal{G}$ (\textbf{Dependency Graph}): The final directed graph $G = (V, E)$. The vertices $V$ are the components extracted from the code and the edges $E$ are the semantic dependency relationships (inheritance, composition, usage).
\end{itemize}

\subsection{Analysis Phases}
\begin{itemize}
    \item $\Phi$ (\textbf{Full Analysis}): The overall mathematical function that defines the architectural pipeline ($\Phi = \gamma \circ \rho_2 \circ \rho_1 \circ \alpha \circ \varepsilon$).
    \item $\varepsilon$ (\textbf{Extraction}): The language-dependent phase that translates the CST of a single file into modules $\mathcal{M}$.
    \item $\alpha$ (\textbf{Aggregation}): The function that merges the modules extracted from individual files into a single global ecosystem.
    \item $\rho_1$ (\textbf{Symbol Definition}): Phase 1 of resolution, which traverses the unified IR to build the static scope tree $\mathcal{T}$, registering all symbols (types, modules, aliases) and base inheritance relationships (e.g., Rust's \texttt{Deref} trait).
    \item $\rho_2$ (\textbf{Name Resolution}): Phase 2 of resolution, which leverages the context $\Gamma$ (based on $\mathcal{T}$) to algebraically evaluate each query $\mathcal{Q}$ and resolve the type references $\tau$, producing validated absolute paths.
    \item $\gamma$ (\textbf{Graphing}): The final phase that converts the unified and resolved object model into the dependency graph $\mathcal{G}$.
\end{itemize}
